\section{Unidad 2}

\subsection{¿Qué condiciona el incremento de la corriente de polarización $I_F$?}

Conforme el voltaje de polarización en directa se incrementa, la corriente continúa incrementándose muy rápidamente, aunque el voltaje a través del diodo se incrementa sólo gradualmente por encima de $0.7\,\text{V}$. Este pequeño incremento en el voltaje del diodo por encima del potencial de barrera se debe a la caída de voltaje a través de la resistencia dinámica interna del material semiconductor.

\subsection{¿Cómo se calcula la corriente en directa $I_F$?}

\subsection{¿Cómo se calcula la corriente en directa $I_F$?}

\subsection{¿Cómo se calcula la corriente en directa $I_F$?}

\subsection{¿A qué frecuencias cree que afectarán más al diodo las capacidades parásitas?}

En el diodo los niveles de capacitancia parásita son los que tienen un mayor efecto.

A bajas frecuencias y a niveles relativamente bajos de capacitancia, la reactancia de un capacitor, determinada por:
\[
X_C = \dfrac{1}{2\pi f C}
\]

en general es tan alta que se le puede considerar de magnitud infinita, representada por un circuito abierto e ignorada.

A altas frecuencias, sin embargo, el nivel de $X_C$ puede reducirse al punto de que creará una trayectoria de “puenteo” de baja reactancia. Si esta trayectoria de puenteo ocurre a través del diodo, en esencia puede evitar que éste afecte la respuesta de la red.

\subsection{¿Qué cree que sucede en la polarización en directa?}

En la región de polarización en directa predomina el efecto de capacitancia que depende directamente de la velocidad a la cual se inyecta la carga en las regiones justo fuera de la región de empobrecimiento.

El resultado es que los niveles incrementados de corriente aumentan los niveles de capacitancia de difusión. Sin embargo, los niveles incrementados de corriente reducen el nivel de la resistencia asociada y la constante de tiempo resultante ($\tau = RC$), la cual es muy importante en aplicaciones de alta velocidad, no llega a ser excesiva.

\subsection{¿Qué sucede si queremos pasar de la polarización directa a inversa?}

Si el voltaje aplicado se tiene que invertir para establecer una situación de polarización en inversa, de algún modo nos gustaría ver que el diodo cambia instantáneamente del estado de conducción al de no conducción.

Sin embargo, por el gran número de portadores minoritarios en cada material, la corriente en el diodo se invierte como se muestra en la figura.


